\documentclass[12pt]{article}
\usepackage{graphicx}
\usepackage{amsmath}
\usepackage{url}
\usepackage{booktabs}
\usepackage[margin=1in]{geometry}
\usepackage{natbib}
\usepackage{xurl}

\title{Uneven Legalization, Spatial Demand Spillovers, and the Regulation of Entry:\\
Evidence from U.S.\ Cannabis Markets}
\author{Mahyar Rezazadeh \\ Babak Sahragard}
\date{August 2026}

\begin{document}
\maketitle

\begin{abstract}
\noindent Uneven state-level legalization of recreational cannabis creates a spatial demand spillover: counties
near a still-prohibited neighboring state face captive cross-border demand that interior counties of the same
state do not. We ask two questions about this spillover. First, does it produce a border premium in legal
cannabis markets. Second, does entry regulation determine which margin absorbs that premium -- existing
sellers' revenue, or the number of sellers and the market's total size. Using a county-year panel spanning
Colorado, Washington, Oregon, Montana, and New Mexico (2014--2026), with a border classification that is
time-varying rather than a fixed snapshot -- accounting for New Mexico's 2022 legalization, Oregon's 2015
legalization, and Oregon's own 2017--2018 transition relative to California and Nevada -- we find a robust
border premium on revenue per dispensary (+48\%, county-clustered $p=0.11$), revenue per capita (+193\%,
$p<0.001$), and dispensary density (+101\%, $p<0.001$). Whether entry is capped, using an institutional
classification that treats binding local-government control over licensing as equivalent to a state cap, does
not moderate the revenue-per-store margin, but it sharply moderates the other two: capped entry reduces the
border premium's multiplier on revenue per capita by 54\% in relative terms (from roughly 193\% to roughly
34\% in level terms, $p=0.009$) and on dispensary density by 56\% in relative terms ($p<0.001$). Oregon's own
within-state entry freeze in January 2022 produces the same-signed, if
underpowered, pattern. A dispensary-level panel built for New Mexico -- the one sample state where individual
firm revenue is observable historically -- confirms the border premium is not a county-aggregation artifact
(+141\% at the level of the individual dispensary, $p=0.001$, significant in every sample year). The results
support a specific mechanism -- entry regulation determines whether spatial
demand spillovers are absorbed by incumbent revenue or by market expansion -- while also reporting, without
forcing a cleaner story, that the mechanism does not extend to the intensive margin the way it was originally
hypothesized to, and that an alternative cut isolating specifically local-government-imposed restrictions
preserves the density result under county clustering but not the other two margins.
\end{abstract}

\section{Introduction}
\label{sec:intro}

By the end of 2025, 24 U.S.\ states had legalized recreational cannabis, though Virginia had not yet
established an operating retail market, leaving 23 states with a functioning legal market and 27 states,
plus the remaining territory, where recreational retail sale remained prohibited. This patchwork did not
arrive all at once, and it has not stopped changing: a state that borders only prohibition-states in one year
can border a freshly legal neighbor the next. That timing creates a natural, if underused, source of
identifying variation. A county sitting near a still-prohibited state faces a captive pool of residents who
cannot buy cannabis legally at home; a county in the interior of the same legal state does not. If cannabis
demand is at all price- or access-sensitive across a state line -- which the existing literature on cannabis,
alcohol, and other regulated goods says it is -- that captive demand should show up somewhere in the legal
market's own data.

Where it shows up, however, is not obvious, and that is the question this paper is really about. A dispensary
facing extra demand from across the border can respond in (at least) two ways: it can sell more at roughly the
same price and margin, raising its own revenue, or new dispensaries can open nearby to meet the extra demand,
raising the county's total revenue and store count without necessarily raising revenue at any one store. Which
of these happens is not a fact about demand; it is a fact about supply-side regulation. A state that caps the
number of licenses, whether by an explicit numerical ceiling or by letting local governments block entry
outright, closes off the second channel by construction. A state that does not cap entry leaves it open. This
paper's central hypothesis, following directly from that logic, is that cross-border demand creates a revenue
premium in border counties, and that entry regulation determines whether that premium is absorbed by existing
sellers (revenue per store) or by market expansion (revenue per capita and dispensary count).

We test this using a county-year panel covering five states -- Colorado, Washington, Oregon, Montana, and New
Mexico -- assembled from each state's own revenue-reporting agency (2014--2026, depending on the state's own
retail-launch date), together with an institutional classification of each state's entry-licensing regime built
from the actual statutes and regulations, not from a state's own label for itself. Two features of the design
depart from prior work in this space, including our own earlier work on Colorado and Washington alone. First,
a county's exposure to a nonlegal-state border is treated as time-varying: because several of the sample
states' neighbors themselves legalized cannabis partway through the sample window (New Mexico in 2022, Oregon
in 2015, and, more narrowly, California and Nevada relative to Oregon's own 2017 observations), a border
classification fixed at a single snapshot date misclassifies a meaningful share of county-years. We recompute
it year by year. Second, entry regulation is classified by where the binding constraint on licensing actually
sits, not by whether a state imposes an explicit numerical cap. Colorado has no state-level cap, but roughly
half its counties ban cannabis retail outright through local option, a restriction documented independently by
\citet{PayanBrownSong2021}; Montana has no numerical cap either, but its 2021 enabling statute requires
counties that voted against the state's 2020 legalization initiative to affirmatively opt in before a
dispensary can locate there. Both are coded as capped-entry states for this reason, a coding choice we explain
and defend in Section~\ref{sec:institutional} rather than assert.

We find, first, that the border premium is real and does not depend on regulatory regime: revenue per store,
revenue per capita, and dispensary density are all significantly higher in border counties, pooling across the
five states. Second, entry regulation moderates two of these three margins in exactly the direction the
hypothesis predicts, and does so with more precision than the underlying border effect itself: the interaction
between border exposure and capped entry is a 54\% relative reduction in the border premium's multiplier on
revenue per capita ($p=0.009$) and a 56\% relative reduction on dispensary density ($p<0.001$) -- almost identical magnitudes, which
is itself informative, since revenue per capita is arithmetically close to the product of density and revenue
per store, so a shared suppression of that size on both density and per-capita revenue but not on per-store
revenue is consistent with the premium moving through the quantity margin specifically. Third, and reported
honestly rather than smoothed over, the same interaction on revenue per store is a precise, stable null: the
border premium on revenue per store is essentially identical in capped and free-entry states. The original
hypothesis predicted regulation would push the premium into the per-store margin specifically when entry is
capped; the data show only that regulation blocks the quantity margin, not that it redirects the premium into
price. We take this as a real, if partial, confirmation, and say plainly where the simple version of the story
does not hold.

The rest of the paper proceeds as follows. Section~\ref{sec:lit} positions the paper against the existing
literature on cross-border cannabis demand and on cannabis entry regulation, two literatures that have
developed independently of each other, and gives an honest assessment of what is and is not new here.
Section~\ref{sec:institutional} describes the five sample states' legalization timing and entry-licensing
institutions in detail. Section~\ref{sec:data} describes the data and the construction of the time-varying
border measure. Section~\ref{sec:methods} lays out the empirical strategy. Section~\ref{sec:results} reports
results. Section~\ref{sec:discussion} discusses magnitude, mechanism, and limitations, and
Section~\ref{sec:conclusion} concludes.

\section{Literature}
\label{sec:lit}

Two literatures bear directly on this paper's question, and, to our knowledge, neither has been connected to
the other.

The first documents that cannabis demand crosses state borders when legal status differs. \citet{HaoCowan2020}
show that recreational legalization draws substantial cross-border sales from residents of neighboring
nonlegal states. \citet{HansenMillerWeber2020} show the mechanism operating in the other direction: Washington
retailers near the Oregon border lost roughly 36 percent of their sales when Oregon's own retail market opened
in October 2015, evidence that the earlier premium those counties enjoyed was genuinely a captive-demand effect
and not some fixed, unrelated feature of border geography. \citet{BaggioChongKwon2020} document a related
cross-border effect between cannabis and a different regulated good, showing that counties in states that
legalized medical marijuana reduced alcohol purchases relative to nearby counties across the state line, using
a border-county-pair design in the tradition of \citet{DubeLesterReich2010}. Outside cannabis, cross-border
shopping in response to tax and regulatory differences is a mature empirical literature; \citet{FribergEtAl2024}
is a recent addition, using Norway's COVID-19 border closures as a natural experiment to show that cross-border
shopping measurably erodes domestic commodity tax revenue, evidence that the underlying mechanism -- consumers
respond to a nearby jurisdiction's more permissive rules -- is not specific to cannabis or to the United States.

The second literature studies how entry regulation shapes market structure and rents once a cannabis market
exists, independent of any cross-border demand story. \citet{HollenbeckUetake2021} document that Washington's
capped licensing regime supports meaningful markups, direct evidence that a binding cap changes market conduct,
not just market structure. \citet{Thomas2019LicenseQuotas} builds a structural model of Washington's own license
quota and estimates that free entry would raise total surplus by roughly 18 percent relative to the actual
capped and partly randomized allocation, decomposing the loss into geographic misallocation and the randomness
of the lottery itself. This is, as far as we can establish, the closest existing paper to the entry-regulation
half of our own question, and it studies exactly the license-cap mechanism this paper relies on -- but with no
demand shock to test it against. The general industrial-organization literature on entry and market structure
supplies the theoretical scaffolding for why a cap should matter at all: \citet{BresnahanReiss1990} and
\citet{BresnahanReiss1991} establish that competitive conduct changes most sharply with the arrival of a
market's second or third entrant; \citet{Berry1992}, \citet{Mazzeo2002}, and \citet{Seim2006} extend this to
settings with firm heterogeneity; \citet{ToivanenWaterson2005} shows a rival's existing presence predicts a
firm's own entry in a dynamic setting; \citet{BerryReiss2007} surveys the tradition. None of this literature
has been applied to cannabis border markets specifically, and none of it pairs an entry-regime moderator with an
exogenous demand shock of the kind cross-border legalization timing provides.

The gap, then, is specific and modest, not a new mechanism. The cross-border demand literature has established
that demand moves across a cannabis-policy border; it has not asked where that demand goes once it arrives --
whether it shows up as price, as quantity, or as some combination determined by supply-side rules the demand
literature does not model. The entry-regulation literature has established that caps distort market structure
and rents; it has not had an exogenous demand shock, of the kind neighboring-state legalization timing
naturally supplies, to test that distortion against. This paper's contribution is bridging those two
literatures with a design built for exactly that purpose, not a new theoretical result. We say this directly
because the honest version of a contribution claim is more defensible than an inflated one, and because the
paper's own findings -- a real interaction on two margins, a precise null on a third -- are themselves a case
for measured claims over sweeping ones.

\section{Institutional Background: Entry Regulation in Five States}
\label{sec:institutional}

\input{tables/table_state_summary}

Table~\ref{tab:state_summary} summarizes each sample state's legalization timing, border exposure, data
sources, and entry-regulation regime; this section explains the regime classifications in more institutional
detail, since they are the paper's central moderating variable and rest on judgment calls we want to make
transparent rather than embed silently in a state-year dummy.

\textbf{Colorado} imposes no state-level numerical cap on retail marijuana licenses -- Amendment 64 (2012)
delegates licensing authority to local governments entirely, and the Colorado Retail Marijuana Code confirms no
statewide ceiling exists. But \citet{PayanBrownSong2021}'s county-level census of Colorado and Washington's
local marijuana ordinances finds that, as of April 2019, 48.4 percent of Colorado counties banned all marijuana
facilities outright, and only 23.4 percent allowed every facility type without restriction -- the most
restrictive local-option environment of any state in this sample, notwithstanding the absence of any state cap.
We classify Colorado as capped for its entire panel on this basis: a state law that says ``no cap'' does not
describe a free-entry market if local governments can and routinely do close most of the state's geography to
retail entirely.

\textbf{Washington} caps entry directly. Initiative 502 (2012) set an initial statewide ceiling of 334 retail
licenses; the Washington State Liquor and Cannabis Board raised it to 556 by emergency rule on January 6, 2016,
following a commissioned market-capacity study, and a social-equity program began issuing up to 52 additional
licenses in 2024. The cap level changes three times over the sample window, but the state is never free-entry
at any point in it, so Washington is coded capped throughout.

\textbf{Oregon} is the one state in the sample whose regime genuinely switches mid-panel, and we treat it
separately for that reason (Section~\ref{sec:results}). Retail licensing was open and largely uncapped from the
market's October 2015 launch. The Oregon Liquor and Cannabis Commission stopped accepting new producer
applications in June 2018, and Senate Bill 218 (2019) extended a producer-only moratorium through January 2022;
House Bill 4016, passed in March 2022 with the freeze applied retroactively to January 2, 2022, extended the
freeze to all license types, including retail. House Bill 4121 (March 2024) replaced the blanket freeze with a
permanent numerical formula -- one retail license per 7,500 residents aged 21 and over. We code Oregon free
from its 2015 launch through the end of 2021 and capped from 2022 onward, matching the freeze's own effective
date rather than the later formal cap, since the freeze is what actually stopped new entry.

\textbf{Montana}'s enabling statute, the Marijuana Regulation and Taxation Act (implemented via HB 701, retail
effective January 1, 2022), imposes no numerical ceiling on the number of dispensary locations. But it
restricts new licenses to conversions of existing medical marijuana providers rather than open applications,
and, more importantly for this paper's coding, it requires counties that voted against the state's own
Initiative 190 in November 2020 to affirmatively approve a dispensary locating there before one legally can --
an opt-in requirement functionally equivalent to Colorado's opt-out bans, and, unlike Colorado's, binding from
the very first day of Montana's retail market rather than emerging gradually. A separate, later restriction
(Senate Bill 27, May 2025, a moratorium on new licenses through at least June 2027) further restricts entry, but
the opt-in requirement alone is sufficient to code Montana capped for its entire panel, not merely from 2025.

\textbf{New Mexico} is the cleanest free-entry state in the sample, and the only one where the classification
rests on statutory text rather than a documented empirical local-restriction rate. The Cannabis Regulation Act
(NMSA 1978 \S 26-2C-6(I)) states that the state's Cannabis Control Division ``shall not limit the number of
licensed premises a licensee may occupy or operate'' and ``shall not place any restriction or prohibition on
the number of licensees'' beyond what the Act itself specifies -- an explicit anti-cap provision, not merely
the absence of one. Licensing is administrative rather than competitive: the statute requires the Division to
begin accepting applications for every license type by January 1, 2022, and imposes a 90-day issue-or-deny
window once an application is complete, with no lottery or merit-scoring step. Local governments retain zoning
and time-place-manner authority (NMSA 1978 \S 26-2C-12(A)) but are explicitly barred from banning cannabis
retailers outright: \S 26-2C-12(B)(2) states a local jurisdiction ``shall not... completely prohibit the
operation of a licensee.'' New Mexico is coded free for its entire panel.

\section{Data}
\label{sec:data}

\subsection{Revenue and dispensary-count data}

Colorado's monthly county-level marijuana sales come from the Colorado Department of Revenue's Socrata series
(2014--2025). Washington's come from the Liquor and Cannabis Board's fiscal-year county sales and excise-tax
reports (FY2015--FY2025). Oregon's monthly county sales come from the Oregon Liquor and Cannabis Commission's
own Socrata series (2017--2025); Oregon's revenue-per-store outcome, unlike its revenue-per-capita outcome,
exists only as a 2025 cross-section, since the OLCC's license registry (the only source that tags dispensary
county directly for Oregon) records a license's most recent renewal date rather than its original opening date
and so cannot support a historical panel. Montana's monthly county sales were scraped directly from the
Department of Revenue's static report pages, since no bulk download exists (2022--2025); Montana's
revenue-per-store outcome has the same single-cross-section limitation as Oregon's, for the same underlying
reason. New Mexico's data come from the Cannabis Control Division's Cannabis Reporting Online Portal, which
publishes monthly sales both by city and, uniquely among the five states, by individual licensee -- meaning
New Mexico's revenue-per-store outcome is a genuine historical panel (2022--2026), not a cross-section, built
directly from licensee-level reported sales rather than from a separately constructed dispensary count.
Dispensary counts for Colorado and Washington come from a geocoded license dataset assembled for a companion
analysis, point-in-polygon assigned to counties and counted as open using each state's own reporting-period
convention.

\subsection{Time-varying border classification}
\label{sec:timevarying}

Every prior version of this analysis, including our own, classified a county's border exposure using a single
snapshot of which states were legal as of a fixed reference date. That classification is wrong for any
county-year in which the relevant neighboring state's own legal status differed from its status on the
snapshot date. Three cases in this sample are large enough to matter. Colorado borders New Mexico, which
legalized retail on 2022-04-01; for Colorado county-years before that date, New Mexico is a nonlegal neighbor,
and eight Colorado counties (Costilla, Archuleta, Conejos, Las Animas, Montezuma, La Plata, Rio Grande, and
Alamosa) fall within 70 miles of a New Mexico county, several of which are not within 70 miles of any other
nonlegal state and so would be misclassified as interior under a static post-2022 snapshot. Washington borders
Oregon, which legalized retail on 2015-10-01; fourteen Washington counties, including Clark County (a large,
populous county directly across the Columbia River from Portland), fall within 70 miles of an Oregon county,
and for Washington's earliest fiscal-year observations -- entirely predating Oregon's launch -- these counties
faced genuine captive cross-border demand from Oregon that a static classification omits entirely. Oregon
itself borders California (legal 2018-01-01) and Nevada (legal 2017-07-01); three Oregon counties (Josephine,
Curry, and Jackson, in the state's south) fall within 70 miles of California, bearing on Oregon's 2017
observations specifically, the first year of Oregon's own panel.

We recompute distance to the nearest nonlegal-state county separately for every observation year, using each
neighboring state's own retail-legalization date rather than a single snapshot, and take a county as
border-exposed if that year-specific distance is 70 miles or less. This reclassifies 41 of 1,314 person-capita
county-year observations and 38 of 995 revenue-per-store county-year observations relative to the static
alternative -- not a large share of the sample, but concentrated exactly where it matters most for internal
validity: the early years of Washington's and Colorado's panels and the first year of Oregon's.

\subsection{Entry-regime coding}

Building on Section~\ref{sec:institutional}, we code a binary \texttt{Capped} indicator at the state-year
level: Colorado, Washington, and Montana are capped throughout their panels; Oregon is free through 2021 and
capped from 2022; New Mexico is free throughout. We treat this as the paper's central regulatory
classification and return to its limits -- principally, that it is a state-year rather than a county-year
measure, and so cannot capture Colorado's own within-state variation in local bans directly -- in
Section~\ref{sec:discussion}.

\subsection{Local- versus state-level entry barriers}
\label{sec:localbarrier}

We also construct an alternative classification, \texttt{LocalBarrier}, that isolates specifically \emph{which
level of government} imposes the binding entry restriction, holding fixed the same institutional facts used to
build \texttt{Capped}. Colorado and Montana are coded 1: in both, the restriction that actually binds is a
county-level decision (Colorado's opt-out bans, Montana's opt-in requirement), not a state-set numerical
ceiling. Washington, Oregon, and New Mexico are coded 0: Washington's cap is set directly by state statute with
no local opt-out or opt-in mechanism; Oregon's freeze (Section~\ref{sec:institutional}) is likewise a
state-level action; New Mexico has no restriction at any level. \texttt{LocalBarrier} therefore reclassifies
Washington relative to \texttt{Capped} -- capped under \texttt{Capped}, not local-barrier under
\texttt{LocalBarrier} -- while leaving Colorado's and Montana's classification unchanged. We use it as a
robustness check on whether \texttt{Capped}'s result depends on lumping state-imposed and locally-imposed
restrictions together (Section~\ref{sec:results}).

\subsection{Dispensary-level panel: New Mexico}
\label{sec:nmdispensary}

Every result above uses county-year aggregates: total county revenue divided by dispensary count or
population. To test whether the county-level border premium survives at the level of the individual seller,
rather than being an artifact of how revenue happens to aggregate across dispensaries within a county, we
separately construct a genuine dispensary-level panel for New Mexico -- the only one of the five sample states
where individual-firm revenue is available historically (Section~\ref{sec:data}). New Mexico's Cannabis Control
Division publishes 51 monthly licensee-level files (April 2022 through the most recent available month), each
listing every active licensee's street address and reported medical, adult-use, and total sales. We identify a
dispensary by its normalized street address rather than its licensee name, since the same physical location can
appear under different names in different months (ownership changes, rebrands, or minor name variants), while
different addresses under the same licensee name are genuinely different stores. This yields 1{,}253 unique
dispensary locations. Each dispensary's entry month is the first calendar month its address appears in the
data; 231 dispensaries (18.4 percent) are already present in the panel's first month (April 2022, the
adult-use market's launch date) and are flagged as left-censored -- almost certainly pre-existing
medical-program dispensaries carried into the new adult-use market rather than true new entrants, since their
true entry date is unobserved -- while the remaining 1{,}022 are true new entrants whose opening month we
observe directly. County assignment reuses the same ZIP-to-county crosswalk (the Census Bureau's 2020 ZCTA-to-
county relationship file) already validated for the county-level New Mexico panel, and each dispensary inherits
its home county's existing 70-mile border classification and population from the county universe file used
throughout this paper; we do not build a separate address-level geocoding pipeline for New Mexico, since no
address-level coordinate source exists for this state in our data (unlike Colorado and Washington, where a
separately assembled geocoded license file gives exact coordinates and opening dates, but not store-level
revenue -- Section~\ref{sec:data}). The regression sample below drops a small number of dispensary-years with
zero reported revenue, leaving 1{,}249 dispensaries and 3{,}703 dispensary-year observations. Because New
Mexico's own border exposure does not vary within the sample period (Texas and Arizona remain nonlegal
throughout) and New Mexico's entry regime is free for its entire panel, this exercise cannot test the
entry-regulation interaction; it serves specifically as a firm-level robustness check on the existence of the
border premium itself (Section~\ref{sec:results}).

\section{Empirical Strategy}
\label{sec:methods}

The main specification is
\[
y_{ist} = \alpha + \beta_1 \, \text{Border}_{ist} + \beta_2 \, \text{Capped}_{st} + \beta_3\,
\text{Border}_{ist} \times \text{Capped}_{st} + \gamma \ln(\text{population}_{ist}) + \delta_t + \eta_i + \varepsilon_{ist}
\]
estimated separately for three outcomes: log(revenue per dispensary store), log(revenue per capita), and
log(dispensary count per 1{,}000 residents). $\text{Border}_{ist}$ is the time-varying 70-mile indicator from
Section~\ref{sec:timevarying}; $\text{Capped}_{st}$ is the state-year entry-regime indicator; $\delta_t$ are
year fixed effects throughout. $\beta_3$ is the coefficient of interest: it tests whether the border premium
differs between capped and free-entry regimes.

State fixed effects $\eta_i$ are included for the revenue-per-capita specification, where Oregon's own
within-state regime switch (Section~\ref{sec:institutional}) provides genuine within-state identification of
$\beta_3$ independent of the cross-state contrast. They are omitted for revenue-per-store and dispensary
density: Oregon's per-store data exists only as a single 2025 (capped-period) cross-section
(Section~\ref{sec:data}), so in that sample $\text{Capped}_{st}$ has no within-state variation for any of the
five states and becomes collinear with the state fixed effects, producing an unstable, uninterpretable main
effect. We verified this directly -- including state fixed effects in that specification inflates the
$\text{Capped}$ coefficient to an implausible multiple-order-of-magnitude estimate -- and report the corrected,
lower-dimensional specification rather than the collinear one. This is a real specification constraint imposed
by the data, not a preference.

Standard errors are clustered by county throughout, alongside heteroskedasticity-robust (HC1) standard errors
reported for comparison; with 165--178 county clusters depending on the outcome, clustering is not as thin an
inference problem here as in the two-state version of this design, though we still report both.

We additionally estimate every specification separately by state, without pooling, as a transparency check on
whether the pooled interaction is representative of any single state's own pattern or an artifact of how the
five states average together (Section~\ref{sec:results}). Oregon's own free/capped switch is tested a third
way, entirely within Oregon and independent of the cross-state comparison: we split Oregon's revenue-per-capita
panel at its own January 2022 freeze and estimate the border coefficient separately in each sub-period, then
test the difference formally with a within-Oregon interaction term.

\section{Results}
\label{sec:results}

\subsection{Does a border premium exist?}

\input{tables/table_main_pooled}

Table~\ref{tab:main_pooled} reports the main pooled specification. The border premium is positive and
significant on all three outcomes: revenue per dispensary store is 48 percent higher in border counties (HC1
$p<0.001$; county-clustered $p=0.11$), revenue per capita is 193 percent higher ($p<0.001$ under both), and
dispensary density is 101 percent higher ($p<0.001$ under both). Cross-border demand shows up in this legal
cannabis market, on every margin we can measure it, pooling across five states with very different licensing
regimes, legalization dates, and reporting conventions. This is the paper's first result, and it replicates and
substantially extends the two-state finding from our own earlier work.

\subsection{Does entry regulation determine which margin absorbs it?}

The interaction term is where the entry-regulation hypothesis is actually tested, and the answer is a genuine
mix, not a clean confirmation. On revenue per capita, capped entry reduces the border premium's multiplier by
54 percent in relative terms (from roughly 193 percent to roughly 34 percent in level terms; county-clustered
$p=0.009$) -- the interaction coefficient itself, $-0.781$, converts to a $54.2$ percent reduction in the
$(1+\text{premium})$ multiplier, not a linear subtraction of the two level percentages. On dispensary density,
capped entry reduces the border premium's multiplier by 56 percent in relative terms, essentially the same
magnitude ($p<0.001$). On revenue per store, the interaction is a precise null: 3.5 percent, statistically
indistinguishable from zero ($p=0.90$).

Figure~\ref{fig:premium_by_regime} shows this pattern directly, plotting the border premium separately for
free- and capped-entry state-years on each outcome. The revenue-per-store bars are nearly identical across
regimes; the revenue-per-capita and density bars are not, and the density bar is negative under capped entry --
border counties in capped states have \emph{fewer} dispensaries per capita than their own state's interior, not
merely a smaller premium.

\begin{figure}[htbp]
\centering
\includegraphics[width=0.85\textwidth]{figures/fig1_premium_by_regime.png}
\caption{Border premium by outcome and entry regime. Bars are the border coefficient from a within-regime
regression on each outcome (log-linear, controlling for log population and year fixed effects), converted to
a percentage; error bars are 95\% confidence intervals from heteroskedasticity-robust standard errors.}
\label{fig:premium_by_regime}
\end{figure}

The near-identical size of the density and per-capita interactions is itself informative, not merely
coincidental: revenue per capita is arithmetically close to the product of dispensary density and revenue per
store, so a shared suppression of similar magnitude on density and per-capita revenue, alongside a null on
per-store revenue, is consistent with the entire capped-entry suppression running through the quantity margin
-- fewer stores open in border counties when entry is capped, which mechanically caps per-capita revenue too,
while each store that does exist still earns the same border premium regardless of regime.

\subsection{State-specific transparency check}

\input{tables/table_state_specific}

Table~\ref{tab:state_specific} and Figure~\ref{fig:state_forest} report the same three outcomes estimated
separately within each state, since a pooled interaction can average over considerable heterogeneity. It does
here. Colorado's own density coefficient is negative and significant ($-27$ percent) despite Colorado's large
positive revenue-per-store premium ($+81$ percent) -- exactly the pattern the interaction predicts for a
capped state, county by county rather than only in the pooled average. New Mexico's density coefficient is
strongly positive ($+109$ percent), consistent with free entry. Montana and Washington's own density
coefficients are smaller and, for Montana, not significant at conventional levels, which is unsurprising given
Montana contributes only 23 revenue-per-store observations to the pooled sample; we do not read a null
coefficient from 23 observations as contradicting the pooled result so much as being underpowered to speak to
it either way. Oregon's own confidence intervals are wide on every outcome for the same reason -- it
contributes only a 2025 cross-section to the per-store and density specifications.

\begin{figure}[htbp]
\centering
\includegraphics[width=0.85\textwidth]{figures/fig2_state_forest.png}
\caption{State-specific border coefficients by outcome, with 95\% confidence intervals (HC1). Each point is
from a separate within-state regression.}
\label{fig:state_forest}
\end{figure}

\subsection{Oregon's own regime switch}

\input{tables/table_oregon_beforeafter}

Oregon is the one state whose panel spans a genuine within-state entry-regime switch, and it is the cleanest
possible test of the interaction hypothesis precisely because it does not rely on any cross-state comparison.
Table~\ref{tab:oregon_beforeafter} and Figure~\ref{fig:oregon_beforeafter} split Oregon's own revenue-per-capita
panel at its January 2022 freeze. The border premium is larger and significant in Oregon's free period
(+173 percent, $p=0.007$) and smaller and no longer significant in its capped period (+106 percent, $p=0.10$)
-- the same direction as the pooled cross-state interaction. The formal within-Oregon interaction term itself,
however, is not statistically significant ($p=0.50$ HC1, $p=0.61$ county-clustered), which we read as a power
problem rather than a contradiction: Oregon has only three counties that are ever border-exposed under the
70-mile threshold, and splitting an already-small panel at a single date leaves little room for a precisely
estimated interaction. Oregon's revenue-per-store outcome cannot be compared across the freeze at all, since
Oregon's only per-store data postdates the freeze (Section~\ref{sec:data}); we report this as a real limitation
rather than omit the comparison silently.

\begin{figure}[htbp]
\centering
\includegraphics[width=0.7\textwidth]{figures/fig3_oregon_beforeafter.png}
\caption{Oregon's revenue-per-capita border premium, estimated separately before and after the January 2022
entry freeze, with 95\% confidence intervals.}
\label{fig:oregon_beforeafter}
\end{figure}

\subsection{Robustness: local- versus state-level entry barriers}

\input{tables/table_localbarrier_pooled}

Table~\ref{tab:localbarrier_pooled} re-estimates the same three outcomes' border-by-regime interaction using
\texttt{LocalBarrier} (Section~\ref{sec:localbarrier}) in place of \texttt{Capped}. The density interaction
survives this alternative cut and remains significant under county clustering ($-41.5$ percent, $p=0.037$),
consistent with the main result. The revenue-per-capita and revenue-per-store interactions, however, do not:
both are significant under HC1 standard errors but lose significance once clustered by county (per capita:
$p=0.35$; per store: $p=0.17$). We report this directly as a genuine caveat on the main \texttt{Capped}
classification, not a result to explain away: the paper's headline finding is more robust on the
quantity/density margin specifically than on revenue per capita, and its precise magnitude depends in part on
which government level's restriction is being tested, not only on whether some restriction exists.

Figures~\ref{fig:localbarrier_percapita} and \ref{fig:localbarrier_perstore} show how this comparison evolves
over the sample period, using a pooled year-by-year event-study specification (border $\times$ year
interactions, with log county population, log(1+captive population), and state and year fixed effects,
clustered by county; Oregon excluded from both groups because its per-store and per-capita panels are too
short and irregular to support this design cleanly -- Section~\ref{sec:data}). On revenue per capita, the
local-barrier states (Colorado and Montana) show a substantially larger border premium than the non-local-
barrier states (Washington and New Mexico) through 2019 -- roughly 100 percentage points higher -- but the gap
closes almost entirely by 2024. On revenue per store, the two groups track more closely throughout the sample,
consistent with the pooled per-store null in Table~\ref{tab:localbarrier_pooled}.

\begin{figure}[htbp]
\centering
\includegraphics[width=0.85\textwidth]{figures/fig11_localbarrier_v2_percapita.png}
\caption{Border premium on revenue per capita, by year, local-barrier vs.\ non-local-barrier states (Oregon
excluded). Pooled event-study specification with log county population, log(1+captive population), state and
year fixed effects, county-clustered standard errors.}
\label{fig:localbarrier_percapita}
\end{figure}

\begin{figure}[htbp]
\centering
\includegraphics[width=0.85\textwidth]{figures/fig12_localbarrier_v2_perstore.png}
\caption{Border premium on revenue per store, by year, local-barrier vs.\ non-local-barrier states (Oregon
excluded). Same specification as Figure~\ref{fig:localbarrier_percapita}.}
\label{fig:localbarrier_perstore}
\end{figure}

\subsection{Firm-level robustness: a New Mexico dispensary panel}

\input{tables/table_nm_dispensary_eventstudy}

Every result so far uses county-level aggregates. Table~\ref{tab:nm_dispensary_eventstudy} and
Figure~\ref{fig:nm_dispensary} re-estimate the border premium directly on New Mexico's individual dispensaries
(Section~\ref{sec:nmdispensary}), rather than on county totals, as a check on whether the aggregate result
reflects a genuine seller-level pattern or is an artifact of how revenue happens to sum across dispensaries
within a county. It is genuine: the firm-level border premium is large, positive, and significant in every year
of the panel (113 to 184 percent, all $p<0.02$), and a pooled single-coefficient version of the same
specification recovers $+141.4$ percent ($p=0.001$), close to the county-level New Mexico contribution to the
main pooled estimate. Border dispensaries are also disproportionately represented among true new entrants: 87.2
percent of border dispensaries are observed opening after the panel's start, versus 78.3 percent of interior
dispensaries, consistent with the extensive-margin story under free entry, though this comparison is
descriptive rather than causal, since New Mexico's own entry regime never varies within the sample. Because New
Mexico's border exposure and entry regime are both time-invariant across the panel, this exercise cannot speak
to the entry-regulation interaction that is the paper's central result; its contribution is narrower and more
specific -- confirming that the border premium is not a county-aggregation artifact, in the one state where
individual-firm revenue is actually observable.

\begin{figure}[htbp]
\centering
\includegraphics[width=0.75\textwidth]{figures/fig13_nm_dispensary_eventstudy.png}
\caption{Border premium on individual-dispensary revenue, by year, New Mexico firm-level panel ($n=1{,}249$
dispensaries). County-clustered standard errors; controls for log(county population), log(1+captive
population), and year fixed effects.}
\label{fig:nm_dispensary}
\end{figure}

\subsection{Distance gradient and nonlegal-side population}

\begin{figure}[htbp]
\centering
\includegraphics[width=0.72\textwidth]{figures/fig5_distance_gradient.png}
\caption{Revenue per store against log-distance to the nearest nonlegal-state county, both residualized on
log population, year, and state fixed effects. Points are decile-bin means; the line is the fitted regression
slope.}
\label{fig:distance_gradient}
\end{figure}

Replacing the binary border indicator with continuous log-distance to the nearest nonlegal-state county
recovers the same pattern in a different functional form: revenue per store falls significantly with distance
from the border (slope $-0.265$, $p<0.001$), pooling across all five states (Figure~\ref{fig:distance_gradient}).

\begin{figure}[htbp]
\centering
\includegraphics[width=0.72\textwidth]{figures/fig6_captive_population.png}
\caption{Revenue per store against log nonlegal-side population within 70 miles, residualized on log
population, year, and state fixed effects, pooling all five states.}
\label{fig:captive_population}
\end{figure}

A companion measure -- the total population of nonlegal-state counties within 70 miles of each border county,
rather than distance to the single nearest one -- was the most robust mechanism evidence in our own earlier
two-state analysis of Colorado and Washington, where it survived county-clustered inference when a simple
distance measure did not. Pooling all five states weakens it to a precise null (slope $+0.015$, HC1 $p=0.54$,
county-clustered $p=0.77$; Figure~\ref{fig:captive_population}). We report this reversal directly rather than
retain the earlier, more favorable two-state framing: whatever is generating the border premium in the pooled
five-state sample, it does not scale cleanly with the size of the captive population nearby once New Mexico,
Oregon, and Montana are added to Colorado and Washington.

\section{Discussion}
\label{sec:discussion}

The central, defensible claim from this analysis is narrower than the paper's opening hypothesis, and we think
the narrower claim is the right one to make. Cross-border cannabis demand produces a measurable premium in
legal border-county markets, on every outcome we can observe it on. Entry regulation genuinely moderates that
premium, but specifically on the quantity margin -- how many stores open, and therefore how much revenue a
county earns in total -- not on the price or intensive margin at any single store. The original hypothesis
predicted regulation would redirect the premium from quantity into price when entry is capped; what the data
show is that regulation blocks the quantity response without redirecting anything, and the price-margin premium
simply exists regardless of regime, roughly the same size whether or not a state caps entry. This is a smaller
claim, but it is the one the data support, and we would rather report it than force the larger one.

Three limitations deserve explicit statement. First, the entry-regime variable is coded at the state-year
level, not the county-year level, even though the institutional detail motivating Colorado's and Montana's
capped classification is itself local (county-by-county opt-out and opt-in decisions). A county-level
restriction measure -- the share of a state's border counties that have banned or restricted retail locally --
would let the interaction test operate on genuinely local variation rather than a state-wide average, and is a
natural next step we did not build for this paper. Second, Oregon's own before/after test, the cleanest single
piece of evidence available because it needs no cross-state comparison, is directionally consistent with the
pooled result but not independently significant; readers should weight it as corroborating, not confirmatory,
evidence. Third, the captive-population result reverses sign of statistical significance once the sample
expands from two states to five, a reminder that mechanism evidence built on a small sample does not
automatically generalize, and that we would be wrong to keep citing our own earlier, more favorable version of
that result without flagging the reversal. Fourth, an alternative cut of the same entry-regime classification
-- isolating specifically local-government-imposed restrictions (Section~\ref{sec:localbarrier}) -- preserves
the density result under county clustering but not the revenue-per-capita or revenue-per-store interactions
(Table~\ref{tab:localbarrier_pooled}), a reminder that the interaction's precise magnitude, though not its
overall direction on the quantity margin, is sensitive to exactly how the capped/free distinction is drawn. Set
against these limitations, the New Mexico dispensary-level panel (Section~\ref{sec:nmdispensary}) is a genuine
strength rather than a caveat: it confirms the border premium itself, independent of the entry-regulation
question, is not an artifact of county-level aggregation.

\section{Conclusion}
\label{sec:conclusion}

Uneven legalization creates real spatial demand spillovers in legal cannabis markets: counties near a still-
prohibited neighboring state earn measurably more revenue, on a per-store basis, a per-capita basis, and a
dispensary-count basis, than interior counties of the same legal state. Entry regulation -- coded from the
actual statutes and local-government authority governing each state's market, not from a state's own
description of itself -- determines whether that premium can flow into market expansion: capped-entry states
see a substantially smaller premium on revenue per capita and on dispensary density, while free-entry states'
extra demand shows up as more stores serving more of the market. The premium on revenue at any single store,
however, does not depend on regulatory regime at all, a real departure from the paper's own starting
hypothesis and one we report rather than paper over. The result connects two literatures -- cross-border
cannabis demand, and cannabis entry regulation -- that have so far developed independently of one another, and
does so with a design built specifically to test where captured demand goes once supply-side rules determine
which margin can respond to it.

\bibliographystyle{apalike}
\bibliography{references}

\end{document}
